\chapter{Podstawy teoretyczne i przegląd literatury}
\label{chap:background}
\section{Rozproszone systemy IoT i przetwarzanie brzegowe}
\writingnote{Wyjaśnij role węzła, bramy i stacji treningowej. Porównaj przetwarzanie lokalne i chmurowe, podając źródła oraz kompromisy.}

\section{Komunikacja MQTT}
MQTT opiera komunikację na modelu publish/subscribe. Specyfikacja opisuje między innymi poziomy QoS i mechanizm Last Will~\citep{mqtt2014}.
\writingnote{Omów wybrane mechanizmy protokołu w kontekście telemetrii, duplikatów, statusu urządzeń i uprawnień do tematów.}

\section{Anomalie operacyjne i zdarzenia związane z bezpieczeństwem}
\writingnote{Zdefiniuj anomalie punktowe, kontekstowe i zbiorowe. Oddziel obserwację nietypowego zachowania od ustalenia jego przyczyny.}

\section{Reguły i metody uczenia maszynowego}
Biblioteka scikit-learn udostępnia narzędzia do uczenia i oceny modeli~\citep{scikitlearn2011}. Jej dokumentacja rozróżnia detekcję obserwacji odstających i detekcję nowych nietypowych obserwacji~\citep{sklearnoutliers}.
\writingnote{Wyjaśnij progi i reguły czasowe oraz wybrane 2--3 metody ML. Dla algorytmów odszukaj i przeczytaj publikacje źródłowe; dokumentacja biblioteki nie zastępuje przeglądu badań.}

\section{Stan badań i luka badawcza}
\writingnote{Porównaj publikacje pod względem danych, urządzeń, metod, metryk i ograniczeń. Uzasadnij własny eksperyment bez deklarowania nowości, której nie potwierdza przegląd literatury.}
